\section{Gestore delle ottimizzazioni}

Il componente \textbf{gestore delle ottimizzazioni} riceve le richieste tramite \textbf{SQL}, che è un linguaggio \textit{dichiarativo}. Questo componente trasforma la richiesta in un \textbf{piano di esecuzione}, rappresentato tramite un linguaggio \textit{procedurale}.

Come linguaggio dichiarativo abbiamo visto \textbf{SQL} e il \textbf{calcolo relazionale}; come linguaggio procedurale abbiamo visto l'\textbf{algebra relazionale}.

In questo caso il DBMS determina il miglior piano di esecuzione, ottimizzando rispetto alle \textit{caratteristiche fisiche} del database.

\[
\text{Query}
\;\Rightarrow\;
\text{Analisi lessicale e sintattica}
\;\Rightarrow\;
\text{Ottimizzazione algebrica}
\]

Le ottimizzazioni algebriche comprendono:
\begin{itemize}
	\item \textbf{Anticipo delle selezioni}
	\item \textbf{Anticipo delle proiezioni}
\end{itemize}

Questa fase è \textit{indipendente} dal contenuto della base di dati.

Il gestore delle ottimizzazioni esegue anche un'ottimizzazione che dipende dai \textbf{metodi di accesso}, ovvero dalle strutture fisiche e dalle \textbf{statistiche} (ad esempio, quante righe sono presenti in ogni tabella). Questa fase deve interagire con la parte del DBMS che contiene i \textbf{profili delle tabelle}. Solo dopo aver completato questa ottimizzazione è possibile generare il \textbf{piano di esecuzione}.

\bigskip
\noindent Le operazioni tipiche di accesso supportate dai DBMS sono:
\begin{itemize}
	\item Scansione sequenziale delle tuple di una tabella
	\item Ordinamento delle tuple
	\item Accesso diretto tramite indice
\end{itemize}

\subsection{Scansione sequenziale delle tuple di una tabella}

È la prima operazione che viene eseguita e spesso viene effettuata contestualmente ad altre operazioni:

\begin{itemize}
	\item \textbf{Scan + proiezione} senza eliminazione dei duplicati
	\item \textbf{Scan + selezione} rispetto a un predicato semplice
\end{itemize}

Il costo della scansione sequenziale è definito come:
\[
\text{Costo} = NP(R)
\]
dove:
\begin{itemize}
	\item $NP$ = numero di pagine/blocchi
	\item $NP(R)$ = numero di pagine della relazione/tabella $R$
\end{itemize}

\subsection{Operazione di ordinamento}

L'ordinamento viene utilizzato nei seguenti contesti SQL:

\begin{itemize}
	\item \texttt{ORDER BY} — ordinare il risultato di un'interrogazione
	\item \texttt{DISTINCT} — eliminare i duplicati
	\item \texttt{GROUP BY} — raggruppare le tuple
\end{itemize}

L'ordinamento viene eseguito in memoria secondaria tramite l'algoritmo \textbf{Sort-Merge}, composto da 2 fasi:

\begin{enumerate}
	\item \textbf{Sort interno:} si legge \textit{una pagina alla volta}, caricando un blocco di memoria centrale alla volta. Tutte le tuple della pagina vengono ordinate tramite \textbf{QuickSort} o un algoritmo equivalente. La pagina ordinata risultante viene chiamata \textbf{RUN} e viene salvata in un file temporaneo in memoria secondaria.
	
	\item \textbf{Merge:} si applicano uno o più passi di \textbf{fusione} con cui le RUN vengono unite progressivamente fino a produrre un'\textit{unica RUN ordinata}.
\end{enumerate}

\subsubsection{Esempio di Sort-Merge a 2 vie ($z = 2$)}

\textbf{Parametri:}
\begin{itemize}
	\item $NP = 4$ — numero di pagine della tabella da ordinare
	\item $NB = 3$ — numero di buffer di memoria centrale disponibili
	\item Sort-Merge a \textbf{2 vie} ($z = 2$): ad ogni passo di fusione si uniscono \textbf{2 RUN} alla volta
\end{itemize}

\bigskip
\noindent La tabella da ordinare rispetto alla \textit{prima colonna} è la seguente:

\[
\begin{array}{|c|c|}
	\hline
	g & 24 \\
	a & 19 \\
	d & 21 \\
	\hline
	c & 23 \\
	b & 14 \\
	e & 16 \\
	\hline
	r & 16 \\
	d & 31 \\
	m & 3  \\
	\hline
	p & 2  \\
	d & 17 \\
	a & 14 \\
	\hline
\end{array}
\]

\bigskip
\noindent\textbf{Fase 1 — Sort interno:} si carica e ordina una pagina alla volta. Ogni pagina ordinata costituisce una \textbf{RUN}, salvata in memoria secondaria.

\begin{itemize}
	\item \textbf{RUN 1:} $a:19,\; d:21,\; g:24$
	\item \textbf{RUN 2:} $b:14,\; c:23,\; e:16$
	\item \textbf{RUN 3:} $d:31,\; m:3,\; r:16$
	\item \textbf{RUN 4:} $a:14,\; d:17,\; p:2$
\end{itemize}

\bigskip
\noindent\textbf{Fase 2 — Merge:} si applicano passi di fusione a 2 vie fino a ottenere un'unica RUN ordinata. Il buffer di output ha capacità di una pagina (3 tuple): quando si riempie viene \textbf{scaricato in memoria secondaria} e svuotato.

\bigskip
\noindent\textbf{Passo 1 — fusione a coppie:}
\begin{itemize}
	\item \textbf{RUN 1 + RUN 2} $\longrightarrow$
	$\underbrace{a:19,\; b:14,\; c:23}_{P_2},\;
	\underbrace{d:21,\; e:16,\; g:24}_{P_3}$
	
	\item \textbf{RUN 3 + RUN 4} $\longrightarrow$
	$\underbrace{a:14,\; d:17,\; d:31}_{P_2'},\;
	\underbrace{m:3,\; p:2,\; r:16}_{P_3'}$
\end{itemize}

\bigskip
\noindent\textbf{Passo 2 — fusione finale:}
\begin{itemize}
	\item \textbf{(RUN 1+2) + (RUN 3+4)} $\longrightarrow$
	$a:14,\; a:19,\; b:14,\; c:23,\; d:17,\; d:21,\; d:31,\; e:16,\; g:24,\; m:3,\; p:2,\; r:16$
\end{itemize}

\subsubsection{Costo del $z$-way Sort-Merge}

Il costo è misurato in termini di \textbf{accessi alla memoria secondaria} (letture e scritture). Consideriamo il caso $z = 2$, $NB = 3$.

\bigskip
\noindent Gli accessi possono essere in \textbf{lettura} o \textbf{scrittura}.

\begin{itemize}
	\item \textbf{Fase di sort interno:} si leggono $NP(R)$ pagine e si riscrivono $NP(R)$ pagine ordinate in memoria secondaria.
	
	\item \textbf{Fase di merge:} ad ogni passo si leggono le $NP(R)$ pagine e si riscrivono ordinate. Il numero di passi di merge necessari è:
	\[
	\lceil \log_2 NP(R) \rceil
	\]
	poiché ad ogni passo il numero di RUN da unire si \textbf{dimezza} (nel caso $z = 2$).
\end{itemize}

\bigskip
\noindent Il \textbf{costo complessivo} del $2$-way Sort-Merge è quindi:

\[
\text{Costo} = 2 \cdot NP(R) \cdot (\lceil \log_2 NP(R) \rceil + 1)
\]

dove:
\begin{itemize}
	\item il fattore $2$ tiene conto di lettura \textbf{e} scrittura
	\item $\lceil \log_2 NP(R) \rceil$ è il numero di passi di merge
	\item $+1$ rappresenta la fase di sort interno
\end{itemize}

\subsection{Accesso diretto tramite indice}

\begin{itemize}
	\item \textbf{Condizione atomica di uguaglianza} $A = v$: richiede la presenza di un \textbf{B+Tree} o \textbf{Hash}. Il costo dipende dal tipo di indice utilizzato.
	
	\item \textbf{Selezione con condizione di range} $A \geq v_1 \;\texttt{AND}\; A \leq v_2$: si utilizza un \textbf{B+Tree}.
	
	\item \textbf{Selezione con congiunzione (\texttt{AND}) di condizioni di uguaglianza} $A = v_1 \;\texttt{AND}\; B = v_2$: viene scelto l'indice in base alla \textbf{condizione più selettiva}. Il DBMS determina autonomamente se utilizzare l'indice su $A$ o su $B$ a seconda di quale condizione riduce maggiormente il numero di tuple restituite.
	
	\item \textbf{Selezione con disgiunzione (\texttt{OR}) di condizioni di uguaglianza} $A = v_1 \;\texttt{OR}\; B = v_2$: si possono utilizzare \textbf{più indici in parallelo}, eseguendo successivamente il \textbf{merge dei risultati}.
\end{itemize}
\subsection{Calcolo del costo di un Join}

Il Join è l'operazione più gravosa da eseguire per un DBMS. Gli algoritmi disponibili sono:

\begin{itemize}
	\item \textbf{Nested Loop Join}
	\item \textbf{Merge Scan Join}
	\item \textbf{Hash Based Join}
\end{itemize}

\bigskip
\noindent\textbf{Osservazione:} dal punto di vista \textit{logico} il Join è \textbf{commutativo}, ma dal punto di vista \textit{fisico} il Join \textbf{non è commutativo}: invertire l'ordine degli operandi non cambia il risultato, ma può cambiare significativamente i \textbf{costi di esecuzione}.

\subsubsection{Nested Loop Join}

Date due relazioni $R$ e $S$ su cui è definito il predicato di Join $PJ$, si distinguono:
\begin{itemize}
	\item $R$ — \textbf{relazione esterna}
	\item $S$ — \textbf{relazione interna}
\end{itemize}

\bigskip
\noindent L'algoritmo è il seguente:

\begin{verbatim}
	per ogni tupla t_R di R {
		per ogni tupla t_S di S {
			if (t_R, t_S) soddisfa PJ {
				aggiungi (t_R, t_S) al risultato
			}
		}
	}
\end{verbatim}

\bigskip
\noindent\textbf{Esempio:} siano date le seguenti relazioni con predicato di Join $PJ: R.A = S.A$.

\[
R:
\begin{array}{|c|c|}
	\hline
	A & B \\
	\hline
	22 & a \\
	87 & s \\
	\hline
	45 & h \\
	32 & b \\
	\hline
\end{array}
\qquad\qquad
S:
\begin{array}{|c|c|c|}
	\hline
	A & C & D \\
	\hline
	22 & z & 8 \\
	45 & k & 4 \\
	\hline
	22 & s & 7 \\
	87 & s & 9 \\
	\hline
	32 & c & 3 \\
	45 & h & 5 \\
	\hline
	32 & g & 6 \\
	\hline
\end{array}
\]

\noindent dove $R$ è composta da 2 pagine ($P_1, P_2$) e $S$ da 4 pagine ($P_1, P_2, P_3, P_4$).

\bigskip
\noindent L'algoritmo carica la prima tupla di $R$ e scansiona \textit{per intero} $S$ verificando il predicato; poi passa alla tupla successiva di $R$ e ripete la scansione completa di $S$, e così via. Il risultato è:

\[
\begin{array}{|c|c|c|c|}
	\hline
	A & B & C & D \\
	\hline
	22 & a & z & 8 \\
	22 & a & s & 7 \\
	\hline
\end{array}
\]
E poi cosi via per ogni elemento
\paragraph{Costo del Nested Loop Join}

Per calcolare il costo è necessario conoscere la \textbf{dimensione del buffer} disponibile.

\bigskip
\noindent\textbf{Ipotesi semplificativa:} 1 pagina per caricare $R$ e 1 pagina per caricare $S$.

\bigskip
\noindent Il costo si compone di due termini:
\begin{itemize}
	\item Si legge $R$ completamente \textbf{una volta}: costo $NP(R)$
	\item Per ogni riga di $R$ si legge $S$ completamente: costo $NR(R) \cdot NP(S)$
\end{itemize}

dove $NR(R)$ indica il numero di righe di $R$.

\bigskip
\noindent Il \textbf{costo complessivo} è quindi:

\[
\text{Costo} = NP(R) + NR(R) \cdot NP(S)
\]
\paragraph{Nested Loop Join con Indice}

Se è disponibile un \textbf{indice} costruito sugli attributi di Join di $S$, la scansione completa di $S$ può essere sostituita da un \textbf{accesso tramite indice}.

\bigskip
\noindent Per ogni tupla di $R$, invece di scansionare interamente $S$, si accede all'indice su $S.A$ per recuperare \textit{solo i blocchi di $S$ che contengono tuple con $S.A = R.A$}.

\bigskip
\noindent\textbf{Costo del Nested Loop Join con indice B+Tree su $S.A$:}

\[
\text{Costo} = NP(R) + NR(R) \cdot \left(\text{Prof\_Indice} + \frac{NR(S)}{VAL(A,S)}\right)
\]

dove:
\begin{itemize}
	\item $NP(R)$ — costo invariato: $R$ deve essere caricata completamente
	\item $NR(R)$ — invariato: numero di volte in cui si esegue il ciclo esterno
	\item $\text{Prof\_Indice}$ — profondità del B+Tree: costo di navigazione dell'indice per ogni tupla di $R$
	\item $\dfrac{NR(S)}{VAL(A,S)}$ — \textbf{selettività} dell'attributo $A$ in $S$: stima del numero di pagine di $S$ da caricare per ogni tupla di $R$
\end{itemize}

\bigskip
\noindent La \textbf{selettività} dell'attributo $A$ nella relazione $S$ è definita come:
\[
\text{Selettività}(A, S) = \frac{NR(S)}{VAL(A,S)}
\]
dove $VAL(A,S)$ è il numero di \textbf{valori distinti} di $A$ che compaiono in $S$.
\bigskip
\noindent\textbf{Esempio:} sia $A = \text{NomeRegione}$, con:
\begin{itemize}
	\item $VAL(A,S) = 21$ — numero di regioni italiane (valori distinti di $A$)
	\item $NR(S) = 21.000$ — numero di righe della relazione $S$
\end{itemize}

\noindent La selettività risulta:
\[
\frac{NR(S)}{VAL(A,S)} = \frac{21.000}{21} = 1.000 \text{ tuple per ogni valore di } A
\]

\noindent Assumendo una \textbf{distribuzione uniforme} dei dati e nel \textbf{caso pessimo} (1 tupla per pagina), si devono caricare \textbf{1.000 pagine} di $S$ per ogni tupla di $R$.
\bigskip
\noindent\textbf{Esempio:} siano date le relazioni:
\begin{itemize}
	\item \texttt{RICETTA} — $NP = 12$, $NR = 260$
	\item \texttt{COMPOSIZIONE} — $NP = 200$, $NR = 13.000$
\end{itemize}

\noindent con predicato di Join $PJ: \texttt{RICETTA.CODICE} = \texttt{COMPOSIZIONE.RICETTA}$.

\bigskip
\noindent\textbf{Metodo 1 — Nested Loop Join senza indice}\\
(\texttt{RICETTA} = relazione esterna, \texttt{COMPOSIZIONE} = relazione interna):
\[
NP(\text{Ricetta}) + NR(\text{Ricetta}) \times NP(\text{Composizione})
= 12 + 260 \times 200 = \mathbf{52.012}
\]

\bigskip
\noindent\textbf{Metodo 2 — Nested Loop Join con indice B+Tree su \texttt{Composizione.Ricetta}}\\
(profondità indice $= 2$, $VAL(\text{Ricetta}, \text{Composizione}) = 260$):
\[
NP(\text{Ricetta}) + NR(\text{Ricetta}) \times \left(\text{Prof\_Indice} +
\frac{NR(\text{Composizione})}{VAL(\text{Ricetta},\text{Composizione})}\right)
\]
\[
= 12 + 260 \times \left(2 + \frac{13.000}{260}\right)
= 12 + 260 \times 52 = \mathbf{13.532}
\]

\bigskip
\noindent\textbf{Conclusione:} in questo caso conviene il \textbf{Metodo 2} (Nested Loop Join con indice B+Tree), poiché $13.532 < 52.012$: l'utilizzo dell'indice riduce significativamente il numero di accessi in memoria secondaria.